\documentclass[11pt]{article}

\begin{document}

\title{Manuscript Review}
\author{}
\date{March 8, 2026}
\maketitle

\section{Executive Summary}
This final global pass finds the manuscript's core argument largely coherent from start to finish: it motivates a falsifiable eighth-criterion search, operationalizes candidates with controlled ablations, and interprets outcomes with explicit SESOI framing. The strongest paper-wide quality is the combination of mechanism checks and outcome-level inference, which prevents many common null-result ambiguities.

Remaining problems are mostly cross-section consistency issues rather than foundational design flaws. Top-level claims are occasionally broader than the exact contrasts supporting them; methods-level promises are not always mirrored in results-level reporting; and conclusion language sometimes sounds more definitive than the bounded evidence warrants. These are correctable with tighter scope labels and clearer inferential tiering.

Quantitative coherence is good but can be improved by aligning test families, effect-size interpretation, interval reporting, and multiplicity policy in one consistent narrative. With these refinements, the manuscript can present a strong, transparent contribution that is both methodologically rigorous and textually consistent across Introduction, Methods, Results, Discussion, and Conclusion.

\section{Prioritized Findings}
The findings below are ranked by confidence (highest first).

\begin{enumerate}
\item \textbf{Claim envelope mismatch remains the main paper-wide consistency risk.} \textbf{Confidence: High.}  
Abstract and conclusion claims often read as global null statements, while detailed sections include at least one robustness sham contrast outside SESOI with strong adjusted significance. Keep headline claims explicitly tied to primary enabled-condition contrasts and retain caveats in the same sentence frame.

\item \textbf{Inferential components are individually sound but not fully synchronized across sections.} \textbf{Confidence: High.}  
The manuscript combines Wilcoxon tests, \(d_z\), TOST, and Holm correction, yet the mapping from estimand to decision rule is not always explicit in prose. A single unified reporting schema would remove ambiguity and improve reproducibility.

\item \textbf{Methods-to-results traceability is incomplete for announced secondary analyses.} \textbf{Confidence: High.}  
Mann--Whitney and Cliff's \(\delta\) are described in Methods, but their role is not consistently visible where core conclusions are made. Either summarize them concisely in the main results or narrow the Methods claims.

\item \textbf{Candidate B evidence supports an ecology-conditional null, not a broad class-level dismissal.} \textbf{Confidence: High.}  
The kin-signal observability bottleneck under default cap is a strong mechanistic explanation, and cap=400 helps contextualize it. The paper should consistently frame this as no detected benefit under tested demographic regimes.

\item \textbf{Candidate A mechanism description is interpretable but under-specified for equation-level auditability.} \textbf{Confidence: Medium-High.}  
The EMA state equation is clear, but reproducibility-critical details (parameter bounds, clipping, stability constraints, fixed versus evolved settings) are not all explicit. This weakens independent verification of the ``works-but-not-selected'' interpretation.

\item \textbf{Legitimacy regression conclusions are plausible but sensitive to feature construction choices.} \textbf{Confidence: Medium.}  
High reducibility \(R^2\) values are informative, yet potential proxy overlap with candidate signals may inflate apparent predictability. Nested validation and explicit feature-audit disclosure would strengthen these claims.

\item \textbf{Tone and structural repetition still create avoidable friction.} \textbf{Confidence: Medium.}  
Several numeric claims are repeated across abstract, discussion, and conclusion with similar wording. Tighter compression and clearer separation of confirmatory versus exploratory claims would improve flow and trust.
\end{enumerate}

\section{Section Recommendations}

\subsection*{Abstract and Introduction}
\begin{itemize}
\item Bind each headline null claim to named contrasts in the same clause, and state whether robustness/sham contrasts are in or out of scope.
\item Keep one concise caveat sentence for the cap=400 seasonal sham exception so the summary remains strictly accurate.
\item Use one primary construct label throughout unless a conceptual distinction is explicitly defined.
\end{itemize}

\subsection*{Methods: Equations and Design}
\begin{itemize}
\item For Candidate A, report \(\alpha\) handling, gain/target bounds, and any clipping or saturation in the correction pipeline.
\item Add one explicit correction-term equation so the mapping from EMA state to control action is fully reproducible.
\item For Candidate B, define neighborhood geometry and normalization conventions for \texttt{kin\_fraction} in one compact formal statement.
\end{itemize}

\subsection*{Methods: Inference Framework}
\begin{itemize}
\item Provide a compact table mapping each comparison family to test, effect metric, interval type, and multiplicity correction.
\item Pair TOST decisions with clearly reported 90\% CIs and state adjustment policy for one-sided equivalence tests.
\item Clarify the evidential status of primary, robustness, pre-shock, and exploratory analyses with explicit labels.
\end{itemize}

\subsection*{Results: Candidate A}
\begin{itemize}
\item Add a brief sensitivity summary across plausible memory-parameter regimes to show robustness of the null conclusion.
\item Include one phase-aware seasonal diagnostic to directly test exploitation of predictable periodic structure.
\item Confirm sham calibration quality beyond compute matching, ideally with moment-level signal comparability.
\end{itemize}

\subsection*{Results: Candidate B}
\begin{itemize}
\item Present default-cap and cap=400 findings in parallel structure so readers can compare mechanisms and outcomes directly.
\item If sham/ablated cap=400 contrasts are discussed, present them in the same reporting block as enabled contrasts.
\item Report kin-signal distributions over defined windows, not only endpoint means, to support observability claims.
\end{itemize}

\subsection*{Legitimacy Analyses}
\begin{itemize}
\item Audit proxy sets for mechanical overlap with candidate signals and disclose any feature exclusions or sensitivity outcomes.
\item Use nested cross-validation for model tuning and permutation inference to reduce optimistic performance bias.
\item Add one robustness model with a different link/function class to test dependence on linear-form assumptions.
\end{itemize}

\subsection*{Discussion and Conclusion}
\begin{itemize}
\item Separate confirmatory findings from exploratory diagnostics in distinct bullets before final synthesis claims.
\item Recast strongest claims as bounded inferences tied to tested ecologies, controls, and comparison families.
\item Reduce duplicate numeric recap and reserve detailed numbers for a single canonical table/figure anchor.
\end{itemize}

\subsection*{Global Consistency}
\begin{itemize}
\item Maintain identical comparison scope labels from abstract through conclusion to prevent interpretive drift.
\item Keep terminology, evidential tiering, and caveat placement uniform across sections.
\item Ensure every major claim points to one visible statistical anchor with matching uncertainty language.
\end{itemize}

\end{document}